\chapter{Basis, Dimension \& Matrix Rank}

A \textbf{Basis} is a minimal generating set of linearly independent vectors that spans a vector space. The number of vectors in any basis defines the unique \textbf{Dimension} of the vector space, while \textbf{Matrix Rank} counts the dimension of its row and column spaces.

\section{Basis for a Vector Space}

\begin{definitionbox}{Basis Definition}
A set of vectors $\mathcal{B} = \{\mathbf{v}_1, \mathbf{v}_2, \dots, \mathbf{v}_n\} \subset V$ is a \textbf{Basis} for a vector space $V$ if it satisfies two conditions:
\begin{enumerate}
    \item $\mathcal{B}$ is \textbf{Linearly Independent}.
    \item $\mathcal{B}$ \textbf{Spans} $V$, i.e., $\vspan(\mathcal{B}) = V$.
\end{enumerate}
Every vector $\mathbf{x} \in V$ can be expressed as a \textbf{unique} linear combination of the basis vectors.
\end{definitionbox}

\begin{examplebox}{Standard Basis for $\mathbb{R}^n$}
The standard basis for $\mathbb{R}^3$ is $\mathcal{E} = \{\mathbf{e}_1, \mathbf{e}_2, \mathbf{e}_3\} = \left\{ \begin{bmatrix} 1 \\ 0 \\ 0 \end{bmatrix}, \begin{bmatrix} 0 \\ 1 \\ 0 \end{bmatrix}, \begin{bmatrix} 0 \\ 0 \\ 1 \end{bmatrix} \right\}$.
Any vector $[x, y, z]^T = x \mathbf{e}_1 + y \mathbf{e}_2 + z \mathbf{e}_3$.
\end{examplebox}

\begin{videocard}{IITM BS Lecture: What is a Basis for a Vector Space?}{https://www.youtube.com/watch?v=lhAQIaFOPxA}
Official lecture introducing basis criteria and unique coordinate representations.
\end{videocard}

\begin{videocard}{IITM BS Lecture: Finding Bases for Vector Spaces}{https://www.youtube.com/watch?v=SxPhClO9zSU}
Official lecture demonstrating methods to extract bases from spanning sets.
\end{videocard}

\section{Dimension of a Vector Space}

\begin{theorembox}{Dimension Theorem}
All bases for a given finite-dimensional vector space $V$ contain the \textbf{exact same number of vectors}. This invariant number is called the \textbf{Dimension} of $V$, denoted $\dim(V)$.
\begin{itemize}
    \item $\dim(\mathbb{R}^n) = n$.
    \item $\dim(\mathbb{R}^{m \times n}) = m \cdot n$.
    \item The zero subspace $\{\mathbf{0}\}$ has dimension $0$.
\end{itemize}
\end{theorembox}

\begin{videocard}{IITM BS Lecture: Rank and Dimension for Vector Spaces}{https://www.youtube.com/watch?v=oATXqim4F5Q}
Official lecture defining dimension, subspace dimension bounds ($\dim(W) \le \dim(V)$), and basis extensions.
\end{videocard}

\section{Matrix Rank via Gaussian Elimination}

\begin{formulabox}{Column Rank, Row Rank, and Matrix Rank}
For any matrix $A \in \mathbb{R}^{m \times n}$:
\begin{itemize}
    \item \textbf{Column Rank:} The dimension of the column space $\vspan(\text{columns of } A)$.
    \item \textbf{Row Rank:} The dimension of the row space $\vspan(\text{rows of } A)$.
    \item \textbf{Fundamental Theorem:} $\text{Row Rank}(A) = \text{Column Rank}(A) = \mathbf{\rank(A)}$.
    \item \textbf{REF Calculation:} $\rank(A)$ is equal to the \textbf{number of non-zero rows (or pivot columns)} in the Row Echelon Form (REF) of $A$.
\end{itemize}
\end{formulabox}

\textbf{Data Science Relevance:} In \textbf{Principal Component Analysis (PCA)} and low-rank matrix approximations (SVD), matrix rank tells us the effective number of independent latent signals in a dataset!

\begin{videocard}{IITM BS Lecture: Rank and Dimension using Gaussian Elimination}{https://www.youtube.com/watch?v=0yEKEavVQJE}
Official lecture computing matrix rank by row reducing matrices to REF.
\end{videocard}

\newpage
\section{Practice Exercises}
\textit{Detailed step-by-step solutions are available in the Appendix.}

\begin{enumerate}
\item \textbf{Basis Verification:} Check whether $\mathcal{B} = \left\{ \begin{bmatrix} 1 \\ 1 \end{bmatrix}, \begin{bmatrix} 1 \\ -1 \end{bmatrix} \right\}$ forms a valid basis for $\mathbb{R}^2$. If so, express $\mathbf{v} = \begin{bmatrix} 5 \\ 1 \end{bmatrix}$ as a linear combination of $\mathcal{B}$.

\item \textbf{Finding a Basis for a Subspace:} Find a basis and the dimension for the subspace of $\mathbb{R}^3$:
$$W = \left\{ \begin{bmatrix} x \\ y \\ z \end{bmatrix} \in \mathbb{R}^3 \Big| x + 2y - z = 0 \right\}$$

\item \textbf{Matrix Rank Calculation:} Compute the rank of the $3 \times 4$ matrix $A$ by reducing it to REF:
$$A = \begin{bmatrix} 1 & 2 & 1 & 4 \\ 2 & 4 & 3 & 11 \\ 3 & 6 & 4 & 15 \end{bmatrix}$$

\item \textbf{Full Rank vs Rank Deficient:} A dataset matrix $X \in \mathbb{R}^{100 \times 5}$ has 100 samples and 5 features. What is the maximum possible rank of $X$? If $\rank(X) = 3$, what does this reveal about the features?

\item \textbf{Data Science Application (Collinear Feature Removal):} Given 3 feature vectors:
$$\mathbf{f}_1 = \begin{bmatrix} 1 \\ 2 \\ 3 \end{bmatrix}, \quad \mathbf{f}_2 = \begin{bmatrix} 2 \\ 4 \\ 6 \end{bmatrix}, \quad \mathbf{f}_3 = \begin{bmatrix} 0 \\ 1 \\ 5 \end{bmatrix}$$
Find the rank of the matrix $F = [\mathbf{f}_1 \; \mathbf{f}_2 \; \mathbf{f}_3]$, extract a minimal basis of independent feature vectors, and identify which feature is redundant.
\end{enumerate}